\scanpage{51}
\mspage{4}
Soit \(\Delta\) un diviseur étale de degré \(6\) dans \(X/S\), muni d'une structure octaédrale i.e. d'une \deleted{\unsure{projection} \unsure{finie}} \deleted{de degré \(2\)} \inserted{opération involutive}
\[
 \deleted{\Delta\longrightarrow\Delta'}\quad \underline{a}_{\Delta}:\Delta\longrightarrow\Delta
\]
\deleted{\unsure{où|i.e.} \(\Delta'\) est fini étale de degré \(3\) sur \(S\).} correspondant à une opération libre de \(\{1,\underline{a}_{\Delta}\}\) sur \(\Delta/S\). On dit que cette structure octaédrale est \deleted{\uline{\unsure{homogène}}, \unsure{si}} \inserted{\uline{régulière} dans} \(X\), si le groupe \(\Gamma=\underline{\operatorname{Aut}}_{\rm oct}^{+}(\Delta)\) opère sur \(X\), de façon compatible avec les op. sur \(\Delta\) (cette op de \(\Gamma\) sur \(X\) est alors unique).
% slanted
\margin{Il faut un petit argument pour montrer que tout \unsure{autom|autom.} de la structure oct. \unsure{de} \(\Delta\) induit par un autom de \((X,\Delta)\) est direct : OPS que l'autom fixe les points de \(\{0,\infty\}=\Delta_1\), et on trouve \unsure{alors|dans} les autom de ``\unsure{rotations|rotation}'' par \(\mu_4\), qui sont bien directs\ldots}

\deleted{\unsure{P|Pour}} Pour \(\Delta\) \inserted{diviseur} fixé, on voit qu'il y a au plus une structure octaédrale sur \(\Delta\) qui soit régulière, et il y en a une ssi \(\Delta\) est une ``configuration octaédrale''. Si \(\Delta\) est une telle configuration, on vient de voir \inserted{déjà} qu'elle est munie d'une \deleted{unique} structure octaédrale, régulière (et \uline{tout} autom. de \(X\) qui \deleted{fixe} \inserted{\unsure{invarie|conserve|respecte}} \(\Delta\) \deleted{fixe} \inserted{\unsure{invarie|conserve|respecte}} cette structure octaédrale).
% slanted
\margin{Si on a une configuration octaédrale régulière \(\deleted{\Delta}\) dans un fibré en droites projectives \(X\), alors le revêtement des orientations de \(\Delta\) est can. isom. au revêtement \(\mu_4^{*}\) \deleted{\unsure{de|\(=\)}} \(\operatorname{Spec}\mathcal O_S[i]\unsure{u}/(i^2+1)\) de \(S\)\ldots}
Inversement, si \(\Delta\) est muni d'une str. oct. rég, alors le sous-groupe \(\underline V\) de \(\underline\Gamma=\underline{\operatorname{Aut}}\deleted{\unsure{\rm oct}}_{\rm oct}(\Delta)\) inv. abélien non trivial est un groupe \deleted{loc.} isom. à \(\mathbb D_{2S}\), opérant fidèlement sur \(X\), et on voit aussitôt que \unsure{la|le}
% archive sheet number: 50
